%% Illustris/IllustrisTNG Data, the Converter, and the LaRT AMR Interface
%% Author: Kwang-Il Seon
\documentclass[english,a4paper]{my_memo}
\synctex=1
\usepackage{listings}
\usepackage{xcolor}
\usepackage{booktabs}
\usepackage{multirow}
\usepackage{longtable}

\lstdefinestyle{pythonstyle}{
  basicstyle=\ttfamily\small,
  keywordstyle=\color{blue}\bfseries,
  commentstyle=\color{gray}\itshape,
  stringstyle=\color{olive},
  frame=single,
  breaklines=true,
  numbers=left,
  numberstyle=\tiny\color{gray},
  language=Python,
  morekeywords={self,True,False,None}
}
\lstdefinestyle{shellstyle}{
  basicstyle=\ttfamily\small,
  frame=single,
  breaklines=true,
  language={}
}
\lstdefinestyle{fortranstyle}{
  basicstyle=\ttfamily\small,
  keywordstyle=\color{blue}\bfseries,
  commentstyle=\color{gray}\itshape,
  stringstyle=\color{olive},
  frame=single,
  breaklines=true,
  numbers=left,
  numberstyle=\tiny\color{gray},
  language=[95]Fortran,
  morekeywords={read,write,open,close,allocate,deallocate}
}
\lstset{style=pythonstyle}

\begin{document}

\title{Illustris/IllustrisTNG Data, the Converter, and the LaRT AMR Interface}
\author{Kwang-Il Seon}
\date{2026-05}
\maketitle

\tableofcontents
\bigskip

%=============================================================================
\section{Overview of Illustris and IllustrisTNG}
\label{sec:overview}
%=============================================================================

The Illustris \cite{vogelsberger2014,nelson2015} and IllustrisTNG
\cite{nelson2019,pillepich2018} projects are large-scale cosmological
magneto-hydrodynamic simulations of galaxy formation and evolution.
They use the \textsc{Arepo} moving-mesh code \cite{springel2010}, which
solves the Euler equations on an unstructured Voronoi tessellation that
moves with the local fluid velocity.

The key difference from AMR codes (such as RAMSES) is that the spatial
discretization is an unstructured Voronoi mesh rather than a structured
octree.  Each gas cell is a convex Voronoi polyhedron defined by a
generator point; the cell positions, volumes, and all physical
quantities are stored directly in HDF5 snapshot files.

\subsection{Simulation Variants}

The publicly available simulation suites include:

\begin{center}
\begin{tabular}{llll}
\toprule
\textbf{Suite} & \textbf{Runs} & \textbf{Box (cMpc)} & \textbf{Ref.} \\
\midrule
Illustris    & Illustris-1, -2, -3        & 75         & \cite{vogelsberger2014} \\
IllustrisTNG & TNG50-1 to -4              & 51.7       & \cite{nelson2019} \\
             & TNG100-1 to -3             & 110.7      &  \\
             & TNG300-1 to -3             & 302.6      &  \\
\bottomrule
\end{tabular}
\end{center}

IllustrisTNG improved upon Illustris with updated AGN feedback (kinetic
and thermal modes), refined supernova feedback, and inclusion of
magnetic fields.

%=============================================================================
\section{Data Format}
\label{sec:format}
%=============================================================================

Both Illustris and IllustrisTNG store data as HDF5 files containing
raw snapshots, group catalogs, merger trees, and supplementary data
catalogs.  This section focuses on the snapshot format, which is the
input to the LaRT converter.

\subsection{Snapshot Structure}

Each snapshot is split across multiple HDF5 chunk files.  A single
chunk file has the following top-level groups:

\begin{center}
\begin{tabular}{lp{10cm}}
\toprule
\textbf{HDF5 Group} & \textbf{Contents} \\
\midrule
\texttt{Header}       & Metadata: particle counts, box size, cosmological parameters,
                        redshift, scale factor \\
\texttt{Parameters}   & Simulation setup parameters (optional) \\
\texttt{Configuration}& Code configuration details (optional) \\
\texttt{PartType0}    & Gas cells (Voronoi mesh elements) \\
\texttt{PartType1}    & Dark matter particles \\
\texttt{PartType3}    & Tracer particles \\
\texttt{PartType4}    & Stars and wind particles \\
\texttt{PartType5}    & Black holes \\
\bottomrule
\end{tabular}
\end{center}

API cutout files (downloaded from the TNG web interface) have a
similar structure but contain only particles/cells belonging to the
requested halo or subhalo.

\subsection{Header Attributes}

The \texttt{Header} group contains the following attributes relevant to
the converter:

\begin{center}
\begin{tabular}{llp{7.5cm}}
\toprule
\textbf{Attribute} & \textbf{Type} & \textbf{Description} \\
\midrule
\texttt{BoxSize}        & float64 & Simulation box side length [ckpc/$h$] \\
\texttt{HubbleParam}    & float64 & Hubble parameter $h$ (e.g.\ 0.6774) \\
\texttt{Redshift}       & float64 & Snapshot redshift $z$ \\
\texttt{Time}           & float64 & Scale factor $a = 1/(1+z)$ \\
\texttt{Omega0}         & float64 & Matter density parameter $\Omega_0$ \\
\texttt{OmegaLambda}    & float64 & Dark energy density parameter $\Omega_\Lambda$ \\
\texttt{NumPart\_Total} & int32[6]& Total particle counts per type \\
\texttt{NumPart\_ThisFile} & int32[6] & Particle counts in this chunk \\
\bottomrule
\end{tabular}
\end{center}

\subsection{Gas Cell Fields (\texttt{PartType0})}
\label{sec:fields}

Table~\ref{tab:fields} lists the \texttt{PartType0} datasets used by
the converter, along with their Illustris code units and the conversion
formulas to physical (CGS/kpc/km\,s$^{-1}$) units.

\begin{longtable}{p{4.5cm}p{3.8cm}p{6.2cm}}
\caption{Gas cell fields in \texttt{PartType0} relevant to LaRT.
Comoving coordinates use scale factor $a$ and Hubble parameter $h$
from the snapshot header.}
\label{tab:fields} \\
\toprule
\textbf{Dataset} & \textbf{Code Units} & \textbf{Physical Conversion} \\
\midrule
\endfirsthead
\toprule
\textbf{Dataset} & \textbf{Code Units} & \textbf{Physical Conversion} \\
\midrule
\endhead
\texttt{Coordinates}
  & ckpc/$h$ (comoving)
  & $x_\mathrm{phys}[\mathrm{kpc}] = x_\mathrm{code} \times a / h$ \\[3pt]

\texttt{Density}
  & $(10^{10}\,M_\odot/h) / (\mathrm{ckpc}/h)^3$
  & $\rho_\mathrm{cgs} = \rho_\mathrm{code} \times
     10^{10}\,M_\odot\,h^2 / (a\;\mathrm{kpc}_\mathrm{cm})^3$ \\[3pt]

\texttt{Masses}
  & $10^{10}\,M_\odot/h$
  & $M[\mathrm{g}] = M_\mathrm{code} \times 10^{10}\,M_\odot / h$ \\[3pt]

\texttt{InternalEnergy}
  & (km/s)$^2$
  & $T = (\gamma - 1)\,u\times 10^{10}\,\mu / k_B$
    (see \S\ref{sec:temperature}) \\[3pt]

\texttt{Velocities}
  & km\,$\sqrt{a}$\,/\,s
  & $v_\mathrm{phys}[\mathrm{km/s}] = v_\mathrm{code} \times \sqrt{a}$ \\[3pt]

\texttt{GFM\_Metallicity}
  & dimensionless
  & Total metal mass fraction $Z$ (not relative to solar) \\[3pt]

\texttt{NeutralHydrogen-}\newline\texttt{Abundance}
  & dimensionless
  & Neutral hydrogen fraction $x_\mathrm{HI} \in [0,1]$
    (from TNG on-the-fly chemistry) \\[3pt]

\texttt{ElectronAbundance}
  & dimensionless
  & $x_e = n_e / n_\mathrm{H}$ \\[3pt]

\texttt{Volume}
  & (ckpc/$h$)$^3$
  & $V_\mathrm{phys}[\mathrm{kpc}^3] = V_\mathrm{code} \times (a/h)^3$ \\[3pt]

\texttt{StarFormationRate}
  & $M_\odot$/yr
  & Instantaneous SFR; cells with SFR $> 0$ are on the effective
    equation of state (see \S\ref{sec:ism}) \\
\bottomrule
\end{longtable}

\subsection{Voronoi Mesh Geometry}

Unlike AMR codes where cell geometry is defined by the octree level
and position, Arepo cells are Voronoi polyhedra.  The \texttt{Coordinates}
field gives the Voronoi generator (control point) position, and the
\texttt{Volume} field gives the cell volume directly.  Cell faces and
neighbor connectivity are \emph{not} stored in the snapshot files;
they must be reconstructed from the Delaunay triangulation if needed.

For the purposes of converting to an octree, the key insight is that
the Voronoi tessellation defines a spatial partition: every point
in the simulation domain belongs to exactly one cell, determined by
which generator is closest.  This property is exploited by the
nearest-neighbor assignment in the converter (Section~\ref{sec:assignment}).


%=============================================================================
\section{Unit Conversions}
\label{sec:units}
%=============================================================================

Illustris/TNG stores all quantities in \textbf{comoving coordinates with
$h$-factors}.  The converter transforms them to physical units suitable
for LaRT.

\subsection{Positions}

\begin{equation}
x_\mathrm{phys}[\mathrm{kpc}]
  = x_\mathrm{code}[\mathrm{ckpc}/h] \times \frac{a}{h}
\end{equation}

After conversion, coordinates are recentered so the galaxy center lies
at the origin (centered convention, range $[-L/2, +L/2]$).

\subsection{Density and Hydrogen Number Density}

\begin{align}
\rho_\mathrm{cgs}[\mathrm{g/cm}^3]
  &= \rho_\mathrm{code} \times
     \frac{10^{10}\,M_\odot \times h^2}
          {(a \times \mathrm{kpc}_\mathrm{cm})^3} \\[6pt]
n_\mathrm{H}[\mathrm{cm}^{-3}]
  &= \frac{\rho_\mathrm{cgs} \times X_\mathrm{H}}{m_\mathrm{H}}
\end{align}

where $X_\mathrm{H} = 0.76$ is the primordial hydrogen mass fraction
and $m_\mathrm{H} = 1.6726 \times 10^{-24}$\,g.

\subsection{Temperature}
\label{sec:temperature}

Illustris stores the specific thermal energy
$u$\,[\,(km/s)$^2$\,=\,$10^{10}$\,erg/g\,].  The temperature is
computed from $u$ and the electron abundance $x_e$:

\begin{align}
\mu &= \frac{4\,m_p}{1 + 3\,X_\mathrm{H} + 4\,X_\mathrm{H}\,x_e} \\[6pt]
T[\mathrm{K}] &= \frac{(\gamma - 1)\,u\times 10^{10}\,\mu}{k_B}
\end{align}

where $\gamma = 5/3$, $m_p = 1.6726\times 10^{-24}$\,g, and
$k_B = 1.3807\times 10^{-16}$\,erg\,K$^{-1}$.
When the \texttt{ElectronAbundance} field is unavailable, the converter
assumes fully neutral gas ($x_e = 0$).

\subsection{Velocities}

\begin{equation}
v_\mathrm{phys}[\mathrm{km/s}]
  = v_\mathrm{code}[\mathrm{km}\,\sqrt{a}\,/\,\mathrm{s}]
    \times \sqrt{a}
\end{equation}

\subsection{Volume and Effective Cell Size}

\begin{equation}
V_\mathrm{phys}[\mathrm{kpc}^3]
  = V_\mathrm{code}[(\mathrm{ckpc}/h)^3] \times \left(\frac{a}{h}\right)^3
\end{equation}

The effective spherical radius of a Voronoi cell is:
\begin{equation}
r_\mathrm{eff} = \left(\frac{3\,V_\mathrm{phys}}{4\pi}\right)^{1/3}
\label{eq:reff}
\end{equation}

This is used by the resolution-matching refinement criterion
(\S\ref{sec:refine_resolution}).


%=============================================================================
\section{ISM Treatment: Star-Forming Cells}
\label{sec:ism}
%=============================================================================

In Illustris/TNG, gas cells with $\mathrm{SFR} > 0$ are placed on an
effective equation of state (EOS) that represents the unresolved
multiphase interstellar medium \cite{springel-hernquist2003}.  Their
\texttt{InternalEnergy} encodes an effective pressure, not the actual
ISM temperature.  The computed temperature for these cells can be
$\sim 10^5$--$10^7$\,K, which is unphysical for the ISM.

The converter provides three treatments via \texttt{--sfr-treatment}:

\begin{itemize}
\item \texttt{cap-temperature} (default): cap $T$ at a user-specified
  value (default $2\times 10^4$\,K) for all cells with $\mathrm{SFR}>0$.
  This is the most common approach in Ly$\alpha$ RT post-processing
  \cite{laursen2009}.

\item \texttt{exclude}: remove all star-forming cells.  This is more
  aggressive but avoids any ambiguity.

\item \texttt{as-is}: use the computed temperature without modification.
  Not recommended for Ly$\alpha$ RT but included for completeness.
\end{itemize}

%=============================================================================
\section{Conversion to Octree AMR}
\label{sec:conversion}
%=============================================================================

LaRT uses an octree AMR grid for photon propagation.  The Voronoi mesh
of Illustris is unstructured and cannot be used directly.  The converter
builds an adaptive octree and assigns physical quantities from the
Voronoi data using the following steps.

\subsection{Step 1: Data Loading and Unit Conversion}

The converter reads \texttt{PartType0} fields from the input HDF5 file
(either a snapshot chunk or an API cutout) and converts all quantities
to physical units as described in Section~\ref{sec:units}.

Coordinates are recentered so that the galaxy center is at the origin.
If \texttt{--center} is not specified, the centroid of the gas cell
distribution is used.  If \texttt{--boxsize} is not specified, the
bounding box of the data is computed with a 10\% margin.

\subsection{Step 2: KD-Tree Construction}

A KD-tree \cite{bentley1975} is built from the Voronoi generator
positions using \texttt{scipy.spatial.cKDTree}.  This tree enables
$O(\log N)$ nearest-neighbor queries from any point in the domain.

\subsection{Step 3: Adaptive Octree Refinement}
\label{sec:refine}

The octree is built using the \texttt{AMRGrid} class from
\texttt{AMR\_grid.py}.  Two refinement criteria are available, applied
in sequence:

\subsubsection{Density and Velocity Gradient Refinement}
\label{sec:refine_gradient}

A cell is refined if the relative variation across its sub-cell probe
points exceeds a threshold:
\begin{equation}
\delta = \frac{q_\mathrm{max} - q_\mathrm{min}}
              {q_\mathrm{max} + q_\mathrm{min} + \epsilon}
       > \delta_\mathrm{thresh}
\end{equation}
where $q$ is the density (or velocity magnitude), and $\epsilon$ is a
small floor to prevent division by zero.  The default threshold is
$\delta_\mathrm{thresh} = 0.3$.  The density and velocity criteria are
combined with logical OR: a cell is refined if \emph{either} criterion
is satisfied.

All cells are first uniformly refined to a minimum level
(\texttt{--level-min}, default 3) before the gradient criterion is
applied.

\subsubsection{Resolution-Matching Refinement}
\label{sec:refine_resolution}

If \texttt{--match-resolution} is specified, an octree cell is
additionally refined whenever its half-width exceeds a factor $f$
(default 1.0) times the effective Voronoi cell radius $r_\mathrm{eff}$
(Equation~\ref{eq:reff}) of the nearest Voronoi cell:
\begin{equation}
h_\mathrm{cell} > f \times r_\mathrm{eff,nearest}
\end{equation}

This ensures the octree resolution matches or exceeds the native
Voronoi resolution of the simulation.

\subsection{Step 4: Property Assignment}
\label{sec:assignment}

Physical quantities ($n_\mathrm{H}$, $T$, $v_x$, $v_y$, $v_z$, and any
requested physics columns) are assigned to each grid point (octree
leaf-cell center or Cartesian cell center) by one of two methods,
selected with \texttt{--resample-method}.

\subsubsection*{Nearest-neighbor (default)}

With \texttt{--resample-method nearest} (the default), the value at a
grid point is taken from the single Voronoi generator nearest to it,
found by a KD-tree query.  This is the \textbf{physically exact}
assignment for a Voronoi tessellation: by definition, the Voronoi cell
of generator $\mathbf{g}_i$ is the set of all points closer to
$\mathbf{g}_i$ than to any other generator, so the nearest-neighbor
query recovers the correct cell ownership.  Sharp density and velocity
discontinuities that the moving mesh preserves (contact
discontinuities, shock fronts, outflow boundaries) are reproduced
without blurring.

\subsubsection*{Gaussian kernel smoothing}

With \texttt{--resample-method gaussian}, the value at a grid point
$\mathbf{r}$ is a volume-weighted Gaussian average over the $K$ nearest
Voronoi cells (\texttt{--kernel-neighbors}, default $K=32$):
\begin{equation}
A(\mathbf{r}) =
  \frac{\sum_{i=1}^{K} V_i\,A_i\,
        \exp\!\left(-d_i^2 / 2 h_i^2\right)}
       {\sum_{i=1}^{K} V_i\,
        \exp\!\left(-d_i^2 / 2 h_i^2\right)},
\label{eq:gauss_gather}
\end{equation}
where $d_i = |\mathbf{r} - \mathbf{g}_i|$, $V_i$ is the Voronoi cell
volume, and $h_i$ is the cell-by-cell smoothing length.  The cell volume
weight $V_i$ keeps the average mass-conservative when cells of very
different sizes contribute.  If the weight sum underflows to zero (a
grid point far from all $K$ neighbors), the converter falls back to the
nearest-neighbor value.

The smoothing length is \textbf{adaptive} by default, scaling with the
local Voronoi resolution,
\begin{equation}
h_i = f_\mathrm{kern}\,r_{\mathrm{eff},i},
\label{eq:adaptive_h}
\end{equation}
where $r_{\mathrm{eff},i}$ is the effective cell radius
(Equation~\ref{eq:reff}) and $f_\mathrm{kern}$ is set by
\texttt{--kernel-size-factor} (default 1.0).  Because TNG/Arepo cells
span two to three orders of magnitude in size between the dense ISM and
the diffuse halo, the adaptive length smooths each region on its own
native scale.  A single fixed smoothing length (in output coordinate
units) can instead be forced with \texttt{--kernel-size~$H$}, which
overrides the adaptive scaling and sets $h_i = H$ for all cells.

Gaussian smoothing is appropriate when a continuous field is desired
(e.g., to suppress the cell-to-cell ``staircase'' in coarsely refined
regions, or to mimic SPH-like kernel fields).  It necessarily softens
the sharp Voronoi discontinuities, so nearest-neighbor remains the
default for fidelity to the moving-mesh structure.

\subsection{Step 5: Optional Physics Computation}

When \texttt{--compute-physics} is specified, the following optional
columns are computed and appended to the output:

\begin{itemize}
\item \textbf{Neutral fraction} ($x_\mathrm{HI}$): from TNG native
  \texttt{NeutralHydrogenAbundance} (\texttt{--ionization from\_illustris}),
  or computed via the CIE formula (\texttt{--ionization cie}), or set
  to 1 (\texttt{--ionization full\_neutral}).

\item \textbf{Electron density} ($n_e$): from TNG
  \texttt{ElectronAbundance} (if using \texttt{from\_illustris}), or
  $n_e = n_\mathrm{H}(1 - x_\mathrm{HI})$.

\item \textbf{Emissivity} (controlled by \texttt{--emissivity-line}):
  \begin{itemize}
  \item \texttt{lya} (default): Case~B Ly$\alpha$ recombination plus
    collisional excitation.
  \item \texttt{civ}: CIE C\,\textsc{iv} $\lambda\lambda$1548,1550
    collisional excitation.  Ion fraction from Gnat \& Sternberg (2007),
    collision strength $\Omega_\mathrm{eff} = 7.5$ from Aggarwal \& Keenan
    (2004), C/H solar abundance from Asplund et al.\ (2009).
    Requires metallicity.
  \item \texttt{ovi}: CIE O\,\textsc{vi} $\lambda\lambda$1032,1038
    collisional excitation.  Ion fraction from Gnat \& Sternberg (2007),
    collision strength $\Omega_\mathrm{eff} = 3.8$ from Aggarwal \& Keenan
    (2004), O/H solar abundance from Asplund et al.\ (2009).
    Requires metallicity.
  \end{itemize}
  The CIE metal-line emissivity is:
  $\varepsilon = n_\mathrm{ion} \times n_e \times q(T)$
  \,[photons\,cm$^{-3}$\,s$^{-1}$],
  where $n_\mathrm{ion} = n_\mathrm{H}\,(Z/Z_\odot)\,A_X\,f_\mathrm{ion}(T)$
  and $q(T) = 8.629\times10^{-6}/(g_l\sqrt{T})\;\Omega\;\exp(-\Delta E/kT)$.
  CIE is appropriate for hot gas ($T \gtrsim 10^5$\,K) where these ions
  peak; for cooler gas the UV background dominates and PIE tables would
  be more accurate.

\item \textbf{Dust density}: Laursen et al.\ (2009) \cite{laursen2009}
  prescription: $n_\mathrm{dust} = (Z/Z_\mathrm{ref})\times
  (n_\mathrm{HI} + f_\mathrm{ion}\,n_\mathrm{HII})$
  with $Z_\mathrm{ref} = 0.0134$, $f_\mathrm{ion} = 0.01$.

\item \textbf{Metallicity}: cell-by-cell $Z$ from \texttt{GFM\_Metallicity}.
\end{itemize}

\subsubsection{PIE vs.\ CIE Neutral Fraction}
\label{sec:pie_vs_cie}

The choice of ionization model (\texttt{--ionization}) has a dramatic
effect on the resulting optical depth and therefore on the simulation
run time:

\begin{description}
\item[\texttt{from\_illustris} (PIE):] TNG computes
  \texttt{NeutralHydrogenAbundance} at runtime using
  photoionization equilibrium (PIE) with the Faucher-Gigu\`{e}re et
  al.\ (2009) metagalactic UV/X-ray background.  At low densities
  ($n_\mathrm{H} \lesssim 10^{-2}\;\mathrm{cm^{-3}}$), the UV field
  ionizes hydrogen efficiently and $x_\mathrm{HI}$ is small.  However,
  at higher densities the gas self-shields against the UV background,
  keeping $x_\mathrm{HI} \approx 1$ even when $T \sim 10^4$\,K.
  Star-forming cells (SFR~$> 0$) sit on an effective equation of state
  and their $x_\mathrm{HI}$ may not be physically meaningful;
  the \texttt{--sfr-treatment} option mitigates this by capping $T$
  or excluding those cells.

  Because self-shielded dense gas retains $x_\mathrm{HI} \approx 1$,
  the PIE neutral fraction typically yields very high line-center optical
  depths ($\tau_\mathrm{pole} \sim 10^{7}$--$10^{8}$), making
  Ly$\alpha$ RT simulations expensive.

\item[\texttt{cie} (CIE):] The neutral fraction is computed from the
  collisional ionization equilibrium (CIE) formula, which depends only
  on temperature.  Gas with $T \gtrsim 2\times10^4$\,K is almost
  fully ionized ($x_\mathrm{HI} \ll 1$), regardless of density.
  This drastically reduces the optical depth in the warm/hot halo gas
  and is appropriate when the UV background is negligible (e.g., within
  a galactic wind or at high redshift before reionization).
  CIE is also useful for testing and benchmarking purposes because the
  lower optical depth leads to faster simulations.

\item[\texttt{full\_neutral}:] All hydrogen is assumed neutral
  ($x_\mathrm{HI} = 1$).  This maximizes the optical depth and is
  mainly useful for idealized tests.
\end{description}

In practice, the physically correct choice depends on the environment
being modeled.  For circumgalactic medium (CGM) studies where the UV
background plays a key role, \texttt{from\_illustris} is preferred.
For hot outflows or testing purposes, \texttt{cie} provides a faster
and still physically motivated alternative.

\subsection{Step 6: Output}

The converter writes the standard LaRT generic AMR format (HDF5, FITS,
or plain text), identical to the output of
\texttt{convert\_ramses\_to\_generic.py}.  See the RAMSES data structure
document for the full schema specification.

\begin{center}
\begin{tabular}{lp{9cm}}
\toprule
\textbf{Dataset} & \textbf{Description} \\
\midrule
\texttt{x, y, z}    & Octree leaf-cell center [kpc, centered] \\
\texttt{level}       & AMR refinement level \\
\texttt{gasDen}      & $n_\mathrm{H}$ [cm$^{-3}$] \\
\texttt{T}           & Temperature [K] \\
\texttt{vx, vy, vz}  & Velocity [km/s] \\
\texttt{metallicity} & Metal mass fraction $Z$ (optional) \\
\texttt{xHI}         & Neutral fraction $x_\mathrm{HI}$ (optional) \\
\texttt{n\_e}        & Electron density [cm$^{-3}$] (optional) \\
\texttt{emissivity}  & Line emissivity [cm$^{-3}$\,s$^{-1}$] (optional; Ly$\alpha$, C\,\textsc{iv}, or O\,\textsc{vi}) \\
\texttt{ndust}       & Dust pseudo-number density [cm$^{-3}$] (optional) \\
\bottomrule
\end{tabular}
\end{center}

%=============================================================================
\section{Converter Usage}
\label{sec:usage}
%=============================================================================

The converter script is located at:
\begin{lstlisting}[style=shellstyle]
LaRT_v2.00/python/AMR_grid/convert_illustris_to_generic.py
\end{lstlisting}

\subsection{Basic Usage (Local File)}

\begin{lstlisting}[style=shellstyle]
python convert_illustris_to_generic.py cutout.hdf5 \
    -o galaxy.h5 \
    --output-unit kpc \
    --level-max 8 \
    --compute-physics \
    --ionization from_illustris \
    --plot
\end{lstlisting}

\subsection{API Download Mode}

The converter can download subhalo cutouts from the IllustrisTNG
public API.  A free API key is required (register at
\texttt{tng-project.org/users/register/}).

\begin{lstlisting}[style=shellstyle]
python convert_illustris_to_generic.py \
    --api-key YOUR_API_KEY \
    --simulation TNG50-1 \
    --snap 99 \
    --subhalo-id 0 \
    -o galaxy.h5 \
    --compute-physics
\end{lstlisting}

\subsection{Full Command-Line Reference}

\begin{longtable}{p{4.5cm}p{2.5cm}p{7cm}}
\toprule
\textbf{Flag} & \textbf{Default} & \textbf{Description} \\
\midrule
\endfirsthead
\toprule
\textbf{Flag} & \textbf{Default} & \textbf{Description} \\
\midrule
\endhead

\texttt{input\_file}  & (positional) & Input HDF5 file (cutout or snapshot chunk) \\
\texttt{-o, --output}  & (required) & Output file (.h5, .fits.gz, or .dat) \\
\texttt{--output-unit}  & \texttt{kpc} & Unit of output coordinates \\
\midrule
\multicolumn{3}{l}{\textit{Grid type}} \\
\texttt{--grid-type}  & \texttt{amr} & Output grid type: \texttt{amr} (adaptive octree)
  or \texttt{cartesian} (uniform $N^3$ grid) \\
\texttt{--ngrid}  & 128 & Grid cells per side (only for \texttt{--grid-type cartesian}) \\
\midrule
\multicolumn{3}{l}{\textit{Octree refinement (AMR mode only)}} \\
\texttt{--level-max}  & 8 & Maximum AMR level \\
\texttt{--level-min}  & 3 & Minimum uniform refinement level \\
\texttt{--dens-threshold}  & 0.3 & Density gradient threshold $\delta_\mathrm{thresh}$ \\
\texttt{--vel-threshold}   & 0.3 & Velocity gradient threshold \\
\texttt{--match-resolution}  & off & Enable resolution-matching refinement \\
\texttt{--resolution-factor} & 1.0 & Factor $f$ for resolution matching \\
\midrule
\multicolumn{3}{l}{\textit{Resampling (Voronoi $\rightarrow$ grid value assignment)}} \\
\texttt{--resample-method}
  & \texttt{nearest}
  & \texttt{nearest} (Voronoi-exact) or \texttt{gaussian}
    (volume-weighted kernel smoothing) \\
\texttt{--kernel-size-factor} & 1.0 & Gaussian mode: adaptive
  $h_i = f_\mathrm{kern}\,r_{\mathrm{eff},i}$ \\
\texttt{--kernel-neighbors}  & 32 & Gaussian mode: number of nearest
  cells $K$ gathered per grid point \\
\texttt{--kernel-size}  & auto & Gaussian mode: fixed smoothing length
  $H$ [output units], overriding the adaptive factor \\
\midrule
\multicolumn{3}{l}{\textit{Physics}} \\
\texttt{--compute-physics}  & off & Compute $x_\mathrm{HI}$, $n_e$, emissivity, $n_\mathrm{dust}$ \\
\texttt{--ionization}
  & \texttt{from\_illustris}
  & Source of $x_\mathrm{HI}$: \texttt{from\_illustris}, \texttt{cie},
    \texttt{full\_neutral}, or \texttt{none} \\
\texttt{--emissivity-line}
  & \texttt{lya}
  & Emission line for emissivity column:
    \texttt{lya} (Case~B Ly$\alpha$),
    \texttt{civ} (CIE C\,\textsc{iv} 1548+1550),
    \texttt{ovi} (CIE O\,\textsc{vi} 1032+1038) \\
\midrule
\multicolumn{3}{l}{\textit{ISM treatment}} \\
\texttt{--sfr-treatment}
  & \texttt{cap-\newline temperature}
  & SFR $>0$ cells: \texttt{cap-temperature},
    \texttt{exclude}, or \texttt{as-is} \\
\texttt{--sfr-cap-temp}  & $2\times10^4$ & Temperature cap [K] \\
\midrule
\multicolumn{3}{l}{\textit{Geometry}} \\
\texttt{--center X Y Z}  & auto & Galaxy center [physical kpc] \\
\texttt{--boxsize}  & auto & Box side length [kpc] \\
\midrule
\multicolumn{3}{l}{\textit{API download (optional)}} \\
\texttt{--api-key}    & --- & TNG API key \\
\texttt{--simulation} & TNG50-1 & Simulation name \\
\texttt{--snap}        & 99 & Snapshot number \\
\texttt{--subhalo-id}  & 0  & Subhalo ID \\
\midrule
\multicolumn{3}{l}{\textit{Miscellaneous}} \\
\texttt{--plot}  & off & Generate diagnostic slice plot ($n_\mathrm{H}$
  and $T$): area-filled cross-sections by default \\
\texttt{--plot-scatter}  & off & With \texttt{--plot}, use the old
  scatter-of-cell-centers style (thin $|z|<\mathrm{boxlen}/50$ slab) \\
\bottomrule
\end{longtable}

\subsection{Diagnostic Plot}
\label{sec:diagnostic_plot}

With \texttt{--plot}, the converter writes a two-column PNG comparing
$n_\mathrm{H}$ and $T$ on a central slice.  By default the panels are
\textbf{area-filled cross-sections}: the left column shows the exact
Voronoi field (a fine $512^2$ pixel grid colored by nearest-neighbor
lookup), and the right column shows the resampled grid as it was
written---true octree squares (one rectangle per leaf cell, sized by
its level) in AMR mode, or an \texttt{imshow} of the central plane in
Cartesian mode.  Comparing the two columns reveals how faithfully the
chosen resolution and resampling method reproduce the underlying
Voronoi structure.

Passing \texttt{--plot-scatter} reverts to the older diagnostic style:
a four-panel scatter plot of the cell centers within a thin
$|z| < \mathrm{boxlen}/50$ slab, colored by value.  This is faster for
very large grids but does not fill the cross-sectional area.

\subsection{Suggested LaRT Input}

After conversion, the converter prints a suggested LaRT input file
snippet.  A typical configuration:

\begin{lstlisting}[style=fortranstyle]
&parameters
 par%use_amr_grid    = .true.
 par%amr_type        = 'generic'
 par%amr_file        = 'galaxy.h5'
 par%distance_unit   = 'kpc'
 par%no_photons      = 1e+06
 par%spectral_type   = 'voigt'
 par%source_geometry = 'diffuse_emissivity'
 par%nxfreq          = 121
 par%nxim            = 100
 par%nyim            = 100
/
\end{lstlisting}

With \texttt{par\%source\_geometry = 'diffuse\_emissivity'}, LaRT uses
the \texttt{emissivity} column from the converted file to determine the
photon source distribution.  Alternatively, use \texttt{'point'} for a
central point source.

\subsection{Cartesian Grid Output}
\label{sec:cartesian_grid}

In addition to the adaptive AMR octree, the converter can produce a
uniform $N \times N \times N$ Cartesian grid using
\texttt{--grid-type cartesian --ngrid N}.  The Voronoi cell data is
resampled onto the regular grid using the same
\texttt{--resample-method} as the octree path (nearest-neighbor by
default, or Gaussian kernel smoothing; \S\ref{sec:assignment}).

The output file stores each physical quantity as a separate 3D image
section (HDF5 group with a \texttt{data} dataset, or FITS image
extension), identified by \texttt{EXTNAME}.  Header attributes
(\texttt{BOXLEN}, \texttt{NX}, \texttt{NY}, \texttt{NZ},
\texttt{UNITLCGS}, \texttt{ORIGINX/Y/Z}) are on the first section.

\begin{lstlisting}[style=shellstyle]
python convert_illustris_to_generic.py cutout.hdf5 \
    -o galaxy_cart.h5 --output-unit kpc \
    --grid-type cartesian --ngrid 128 \
    --compute-physics --ionization cie \
    --boxsize 200
\end{lstlisting}

LaRT reads this format via the \texttt{par\%cart\_file} parameter:

\begin{lstlisting}[style=fortranstyle]
&parameters
 par%cart_file       = 'galaxy_cart.h5'
 par%distance_unit   = 'kpc'
 par%source_geometry = 'diffuse_emissivity'
 par%no_photons      = 1e+06
 par%spectral_type   = 'voigt'
 par%nxfreq          = 121
 par%nxim            = 100
 par%nyim            = 100
/
\end{lstlisting}

When \texttt{par\%cart\_file} is set, LaRT uses the standard Cartesian
raytrace (no AMR overhead).  Grid dimensions and box size are read from
the file header, overriding any \texttt{par\%nx}/\texttt{par\%xmax}
values.  Optional columns (\texttt{xHI}, \texttt{emissivity},
\texttt{ndust}) are applied automatically: \texttt{xHI} scales the
H\,\textsc{i} opacity, \texttt{emissivity} populates the diffuse
source, and \texttt{ndust} sets the dust opacity.


%=============================================================================
\section{End-to-End Workflow}
\label{sec:workflow}
%=============================================================================

A typical analysis workflow is:

\begin{enumerate}
\item \textbf{Obtain data}: download a subhalo cutout from the TNG
  public API, or locate a local snapshot chunk file.

\item \textbf{Convert to octree}:
\begin{lstlisting}[style=shellstyle]
python convert_illustris_to_generic.py cutout.hdf5 \
    -o galaxy.h5 --output-unit kpc \
    --level-max 8 --compute-physics \
    --ionization from_illustris \
    --sfr-treatment cap-temperature \
    --boxsize 200 --plot
\end{lstlisting}

\item \textbf{Inspect the grid}: use \texttt{AMRGrid.read} and
  \texttt{slice\_plot} to verify the density and velocity structure.

\item \textbf{Run LaRT}:
\begin{lstlisting}[style=shellstyle]
mpirun -np 64 LaRT.x galaxy_lya.in
\end{lstlisting}

\item \textbf{Analyze output}: use \texttt{read\_lart.py} to read the
  output and produce spectra, surface-brightness maps, etc.
\end{enumerate}


%=============================================================================
\section{Technical Notes}
\label{sec:notes}
%=============================================================================

\subsection{Nearest-Neighbor vs.\ Kernel Smoothing}

For Voronoi data from Arepo, the nearest-neighbor KD-tree query
(\texttt{--resample-method nearest}, default) recovers the exact
Voronoi cell assignment.  Each point in space belongs to exactly one
cell (the one whose generator is closest), so the nearest-neighbor
query returns the correct cell-averaged quantities and preserves the
sharp density and velocity discontinuities the moving mesh is designed
to capture (contact discontinuities, shock fronts, galactic outflow
boundaries).  This is the recommended default for fidelity to the
simulation.

The Gaussian kernel option (\texttt{--resample-method gaussian},
Equation~\ref{eq:gauss_gather}) is provided for cases where a
\emph{continuous} resampled field is preferable to an exact piecewise
one---for example, to suppress the cell-to-cell ``staircase'' that
appears when a steep gradient is resolved by only a few coarse octree
cells, or to compare against SPH-style kernel fields.  Because the
adaptive smoothing length tracks the local cell size
(Equation~\ref{eq:adaptive_h}), each region is smoothed on its own
native scale rather than with a single global length.  Smoothing
inevitably softens the Voronoi discontinuities and slightly lowers
peak densities, and it also reduces the number of refined octree cells
(gradients are gentler), so it trades exactness for smoothness.

\subsection{Memory and Performance}

For a typical TNG cutout with $\sim 10^6$ gas cells:
\begin{itemize}
\item KD-tree construction: $\sim 1$\,GB, $\sim 10$\,s
\item Octree refinement to level 8: $\sim 1$--$5\times 10^6$ leaf cells,
  $\sim 1$--$5$\,min
\item Output HDF5 file: $\sim 50$--$200$\,MB
\end{itemize}

For $\texttt{--level-max} > 9$, the octree can contain $> 10^8$ cells
and may require $> 10$\,GB of memory.  Use with caution.

\subsection{Coordinate Centering}

API cutout files from TNG contain cells from the entire parent halo,
not centered on the subhalo.  The converter automatically recenters
the data by computing the centroid of the gas cell distribution.
For more precise centering (e.g., on the subhalo potential minimum),
use \texttt{--center X Y Z} with the subhalo position from the
group catalog.

\subsection{Comparison with RAMSES Converter}

The Illustris and RAMSES converters produce the same output format and
can be used interchangeably with LaRT.  The key differences are:

\begin{center}
\begin{tabular}{lp{5.5cm}p{5.5cm}}
\toprule
& \textbf{RAMSES} & \textbf{Illustris} \\
\midrule
Input mesh & Octree AMR & Voronoi (unstructured) \\
Conversion & Direct: read leaf cells & Resample: nearest-neighbor or kernel smoothing \\
Fortran converter & Yes & No (Python only) \\
Ionization & CIE only & CIE or TNG native \\
Temperature & Directly from thermal pressure & From InternalEnergy + $x_e$ \\
ISM treatment & Not needed & SFR temperature cap \\
\bottomrule
\end{tabular}
\end{center}


%=============================================================================
\section{Data Access}
\label{sec:access}
%=============================================================================

\subsection{Public API}

Both projects provide REST APIs for querying and downloading data:

\begin{itemize}
\item Illustris: \texttt{https://www.illustris-project.org/api/}
\item IllustrisTNG: \texttt{https://www.tng-project.org/api/}
\end{itemize}

An API key is required (free registration for academic use).  The API
supports searching subhalos by mass, SFR, or other properties, and
downloading cutout HDF5 files for individual halos or subhalos.

\subsection{Direct File Downloads}

Full snapshot files can be downloaded directly from the project
websites.  A single TNG100-1 snapshot at $z=0$ is $\sim 20$\,GB
(split across $\sim 450$ chunk files).

\subsection{Python Helper Scripts}

The official \texttt{illustris\_python} package
(\texttt{https://github.com/illustristng/illustris\_python}) provides
convenience functions for loading subhalos from local snapshot files.
The LaRT converter does \emph{not} depend on this package; it reads
HDF5 files directly with \texttt{h5py}.


%=============================================================================
\begin{thebibliography}{99}
\bibitem{vogelsberger2014}
Vogelsberger, M., Genel, S., Springel, V., et al.\ 2014, MNRAS, 444, 1518
(\textit{Introducing the Illustris simulation: The evolving population of
galaxies across cosmic time}).

\bibitem{nelson2015}
Nelson, D., Pillepich, A., Genel, S., et al.\ 2015, Astron.\ Comput., 13, 12
(\textit{The Illustris simulation: Public data release}).

\bibitem{nelson2019}
Nelson, D., Springel, V., Pillepich, A., et al.\ 2019, Comput.\ Astrophys.\ Cosmol., 6, 2
(\textit{The IllustrisTNG simulations: Public data release}).

\bibitem{pillepich2018}
Pillepich, A., Springel, V., Nelson, D., et al.\ 2018, MNRAS, 473, 4077
(\textit{Simulating galaxy formation with the IllustrisTNG model}).

\bibitem{springel2010}
Springel, V.\ 2010, MNRAS, 401, 791
(\textit{E pur si muove: Galilean-invariant cosmological hydrodynamical
simulations on a moving mesh}).

\bibitem{springel-hernquist2003}
Springel, V., Hernquist, L.\ 2003, MNRAS, 339, 289
(\textit{Cosmological smoothed particle hydrodynamics simulations: A
hybrid multiphase model for star formation}).
Effective EOS model for the multiphase ISM.

\bibitem{laursen2009}
Laursen, P., Sommer-Larsen, J., Andersen, A.\ C.\ 2009, ApJ, 704, 1640
(\textit{Lyman-$\alpha$ radiative transfer with dust}).
Source of the dust prescription and the ISM temperature-cap convention.

\bibitem{bentley1975}
Bentley, J.\ L.\ 1975, Commun.\ ACM, 18, 509
(\textit{Multidimensional binary search trees used for associative searching}).
\end{thebibliography}

\end{document}
